 In a nested trial design, the randomized trial is embedded in a cohort of eligible people who are proposed to participate in the trial, but if they refuse, they are still included in the observational study. 
Two concrete examples of nested designs are the Women Health Initiative~\citep{anderson2003implementation} and the recent study on Medicaid~\citep{degtiar2021conditional}.


In contrast to~\Cref{sec:setting}, the nested design has an extra variable $S \in \{0,1\}$, which is a binary indicator for randomized trial participation.
 More formally, we observe i.i.d. samples from an underlying distribution $\pfullnest$ over $\left(X, U, Y(0), Y(1),Y, S, T\right)$. Further, let $\pnest \defeq \marg(\pfullnest)$ be the marginal distribution  over $(X,Y,S,T)$. We can then  write the marginal distributions over $X$ of the  (restricted) observational study and randomized trial  as 
$$
\prct_X(X) =  \pnest\left(X \mid S=1 \right) \quad \mathrm{and}  \quad 
\pobsres_X(X) = \pnest\left(X \mid S=0, X \in \supp(\prct) \right).$$
This study design has a significant impact on the estimators previously introduced. In particular, the importance weights $w(X)$ can be estimated by pooling the observational study and the randomized trial as follows %Recall  from~\Cref{lemma:identify_ate} that \begin{align*}\mu(\pobsres_X, \pfullobsres)  =\EE_{\prct}\left[ Y\left( \frac{T}{\pi} - \frac{(1-T)}{1-\pi} \right) w(X)\right],
%\end{align*} 
%where the restricted observational study population is defined as 
%$
%\pobsres \defeq \pobs \mid X \in \supp(\prct). 
%$ 
%We show that the importance weights can be identified as follows
\begin{align*}
	w(X) &= \frac{\pobsres(X)}{\prct(X)} \\&= \frac{\pnest(X\mid S=0, X\in \supp(\prct))}{\pnest(X \mid S=1)}\\
	&= \frac{\pnest(S=0\mid X) }{\pnest(S=1\mid X)} \frac{\pnest(S=1 \mid X \in \supp(\prct))}{\pnest(S=0 \mid X \in \supp(\prct))},
\end{align*}
where the sampling probability $\pnest(S\mid X)$ can be identified under a nested study design~\citep{dahabreh2021study}.


